problem:  
Let \(M^n\) be a closed spin manifold with a Riemannian foliation \(\mathcal{F}\) of codimension \(q\) and longitudinal Dirac operator \(D_\mathcal{F}\). Assume every leaf of \(\mathcal{F}\) admits a complete metric of **uniformly positive scalar curvature (PSC)**, but suppose the foliation has **nontrivial holonomy** and the transverse structure is only **leafwise measure‑preserving**, not transversally flat.  

Define the **longitudinal rho class** \(\rho(D_\mathcal{F},g_\mathcal{F})\in K_1(C^*(M,\mathcal{F}))\) as the secondary invariant of the perturbed longitudinal Dirac operator, following Connes–Moscovici.  

**Question:**  
Can nontrivial holonomy cause the longitudinal rho class \(\rho(D_\mathcal{F},g_\mathcal{F})\) to be **nonzero** even when all leaves have uniformly positive scalar curvature?  
Equivalently, is it possible for a foliation with leafwise PSC to carry a *secondary analytic obstruction* coming from nontrivial holonomy, preventing the global vanishing of the longitudinal rho class?  

If such examples exist, how does this secondary obstruction relate to the image of the **Baum–Connes assembly map** for the holonomy groupoid of \(\mathcal{F}\)?  

Why is it a "good" problem:  
This problem builds on the subtle frontier between **positive scalar curvature geometry** and **noncommutative index theory for foliations**, where secondary invariants (rho classes) encode fine information invisible to primary index classes. Classical theorems show that positive scalar curvature kills primary Dirac indices, but the fate of **secondary classes under holonomy and leafwise PSC** is deeply unclear.  

It directly tests whether holonomy can retain analytic curvature information in the \(K\)-theory of the groupoid C\(^*\)-algebra, potentially linking the world of scalar curvature to the **Baum–Connes conjecture** and the structure of groupoid \(K\)-theory. The problem is simple to state—“does holonomy let rho classes survive under PSC?”—yet resolving it would likely require new methods combining analysis on foliations, coarse geometry of groupoids, and higher index theory. It thus has the hallmarks of a *good* problem: a crisp formulation with profound implications connecting geometry, topology, and noncommutative analysis.